% !TEX program = pdflatex
\documentclass[11pt]{article}

% ── Packages ──────────────────────────────────────────────────────────
\usepackage[margin=1in]{geometry}
\usepackage{amsmath,amssymb,amsthm}
\usepackage{booktabs}
\usepackage{graphicx}
\usepackage{hyperref}
\usepackage{cleveref}
\usepackage{xcolor}
\usepackage{caption}
\usepackage{subcaption}
\usepackage{multirow}
\usepackage{enumitem}
\usepackage[numbers]{natbib}
\bibliographystyle{plainnat}

\hypersetup{colorlinks=true,linkcolor=blue!60!black,citecolor=green!50!black,urlcolor=blue!70!black}

% ── Macros ────────────────────────────────────────────────────────────
\newcommand{\Einf}{E_\infty}
\newcommand{\Reval}{R_{\mathrm{eval}}}
\newcommand{\Rtrain}{R_{\mathrm{train}}}
\newcommand{\relLtwo}{\varepsilon_{\mathrm{rel}}}
\DeclareMathOperator*{\argmin}{arg\,min}

% ── Result macros ─────────────────────────────────────────────────────
\newcommand{\totalRuns}{32{,}783}
\newcommand{\mainRuns}{31{,}370}
\newcommand{\ablationRuns}{1{,}413}

% ══════════════════════════════════════════════════════════════════════
\title{%
  Empirical Scaling Laws for Neural Operator Learning:\\
  Architecture, Data, Resolution, and Regime Dependence\\
  on Canonical 1D PDE Tasks%
}
\author{Lucas Perrier}
\date{\today}

\begin{document}
\maketitle

% ══════════════════════════════════════════════════════════════════════
\begin{abstract}
Scaling laws have become central tools for predicting and allocating
resources in deep learning, yet their empirical status in operator
learning---where discretization resolution introduces a third axis
alongside dataset size and model capacity---remains unclear.
We present a controlled empirical study of \totalRuns{} training runs
spanning three canonical 1D PDE benchmarks (Burgers, Darcy flow,
heat diffusion), three neural-operator architectures (FNO, DeepONet,
MLP), and two orders of magnitude along each of dataset size~$N$,
resolution~$R$, and parameter count~$D$.
Every configuration is repeated across nine random seeds, and all
models share an identical training protocol.

Three findings emerge with strong reproducible evidence,
all specific to this 1D setting.
First, \emph{architecture choice dominates all other scaling axes
on these benchmarks}:
FNO wins 143 out of 143 head-to-head comparisons at matched capacity,
often by an order of magnitude---though this dominance is partly
expected given that two of three tasks (Burgers, diffusion) are
periodic and spectrally structured, closely matching FNO's Fourier
inductive bias.
Second, \emph{dataset size is the primary useful lever}:
a descriptive log-linear regression attributes
65--93\% of error variance to~$N$;
capacity contributes modestly and resolution is negligible or harmful.
Third, \emph{higher resolution does not reliably reduce error}:
it actively increases error on smooth problems (diffusion), provides
no signal on elliptic problems (Darcy), and helps only on
advection-dominated problems (Burgers)---a finding confirmed by
a parameter-controlled MLP experiment that isolates resolution from
capacity.
We further show that \emph{scaling-law identification is fragile}:
the AICc-selected functional form shifts with grid density,
and no single law (power, exponential, logarithmic) is robust across
tasks and architectures.
The strongest methodological lesson is that law identification
requires denser grids than are typical in operator-learning studies.
All data, code, and analysis scripts are publicly available.
\end{abstract}


% ══════════════════════════════════════════════════════════════════════
\section{Introduction}
\label{sec:intro}

Neural scaling laws---empirical relationships between performance and
resources such as dataset size, model capacity, and compute---have
reshaped how the machine-learning community allocates training budgets.
The landmark results of \citet{kaplan2020scaling} and
\citet{hoffmann2022training} demonstrated that language-model loss
follows smooth power laws in parameters and tokens, enabling reliable
extrapolation and motivating the ``scaling hypothesis'' that has driven
much of modern AI investment.
Earlier, \citet{hestness2017deep} established similar power-law
relationships across vision, speech, and translation tasks.

Operator learning---the task of learning mappings between function
spaces, typically solution operators of partial differential equations
(PDEs)---differs from these settings in a fundamental way.
In addition to dataset size~$N$ and model capacity~$D$, discretization
resolution~$R$ introduces a \emph{third} scaling axis that has no
direct counterpart in language or vision.
Whether the smooth, universal power laws observed in the token/parameter
plane of language models transfer to the three-axis
$(N, R, D)$ landscape of operator learning is an open empirical
question.

Several recent works have begun to address operator-learning scaling.
\citet{li2021fourier} introduced the Fourier Neural Operator (FNO) and
demonstrated resolution invariance in spectral convolutions.
\citet{lu2021learning} proposed DeepONet based on the universal
approximation theorem for operators~\citep{chen1995universal}.
Yet systematic, controlled comparisons across the full $(N, R, D)$ grid
are scarce: most evaluations fix two axes and vary one, use
architecture-specific hyperparameter tuning, or report results on a
single PDE.
This leaves basic questions unanswered: when does adding resolution
help? How important is model capacity relative to data? Is there a
universal scaling form, or does the relationship depend on the task
and architecture?

This paper presents a large-scale empirical study designed to answer
these questions---and to be honest about what cannot be answered with
current evidence.
We make no claim to a universal scaling law for operator learning.
Instead, we provide \emph{controlled, reproducible measurements} across
a grid of three 1D PDE tasks, three architectures (plus a
parameter-controlled MLP variant), and $\sim$32{,}000 training runs,
all under an identical training protocol.
Every claim is backed by seed-averaged results across nine
random initializations, and we report evidence strength explicitly.

Our main contributions are:
\begin{enumerate}[leftmargin=*,itemsep=2pt]
  \item A \textbf{controlled experimental protocol} covering three
    scaling axes, three tasks, and four model variants with identical
    training configuration and nine-seed replication
    (\Cref{sec:design}).

  \item Evidence that \textbf{architecture choice overwhelms all other
    scaling considerations} for these 1D tasks: FNO outperforms
    DeepONet and MLP in every tested regime (\Cref{sec:architecture}),
    with an explicit discussion of how the periodic, spectrally
    structured nature of these benchmarks favors FNO's Fourier
    inductive bias (\Cref{sec:spectral-bias}).

  \item A demonstration that \textbf{resolution is not a useful scaling
    axis in general}: it helps for Burgers, is noise for Darcy, and
    actively hurts for diffusion---with a parameter-controlled
    experiment confirming the effect is not an artifact of the
    MLP capacity confound (\Cref{sec:resolution,sec:confound}).

  \item Evidence that \textbf{scaling-law identification is unstable
    under sparse grids}: AICc-based model selection across six
    candidate laws reveals no single winner, and the selected law
    shifts with grid density---a methodological finding that cautions
    against over-interpreting any particular functional form
    (\Cref{sec:law-form}).

  \item \textbf{Optimizer and convergence ablations} showing that
    directional findings (resolution helps/hurts, data has bounded
    returns) are stable across Adam, SGD+cosine, and doubled
    early-stopping patience (\Cref{sec:ablations}).

  \item \textbf{Full public release} of all 32{,}783 run logs,
    analysis scripts, and configuration files.
\end{enumerate}

We restrict all claims to the 1D PDE setting studied here.
Generalization to higher dimensions, more complex PDEs, or
larger-scale architectures remains an important open problem.


% ══════════════════════════════════════════════════════════════════════
\section{Related Work}
\label{sec:related}

\paragraph{Neural scaling laws.}
\citet{hestness2017deep} provided the first broad survey of power-law
scaling across deep learning domains. \citet{kaplan2020scaling}
distilled this into clean scaling laws for autoregressive language
models, and \citet{hoffmann2022training} refined the compute-optimal
frontier (``Chinchilla scaling''). These works established the
expectation that test loss follows $E \propto X^{-\alpha}$ in
dataset size, model size, or compute.

\paragraph{Operator learning.}
The FNO~\citep{li2021fourier} learns in spectral space, enabling
resolution invariance via truncation of Fourier modes.
DeepONet~\citep{lu2021learning} uses a branch-trunk decomposition
motivated by the universal approximation theorem for
operators~\citep{chen1995universal}.
Extensions include physics-informed
DeepONets~\citep{wang2021learning}, U-FNO~\citep{wen2022ufno}, and
the General Neural Operator Transformer
(GNOT)~\citep{hao2023gnot}.
These works focus on architectural innovation and accuracy on
individual benchmarks; systematic scaling studies across
$(N, R, D)$ are rare.

\paragraph{Scaling in scientific ML.}
\citet{perrier2025scaling} introduced scaling-law diagnostics for
physics-informed machine learning, proposing multi-law model selection
via AICc. The present paper builds on that diagnostic framework while
shifting the emphasis from methodology to empirical findings: we use
the diagnostic tools to characterize what actually happens when each
axis is scaled, rather than proposing a new fitting methodology.


% ══════════════════════════════════════════════════════════════════════
\section{Experimental Design}
\label{sec:design}

The study is designed to isolate the effect of each scaling axis
($N$, $R$, $D$) while keeping all other factors constant.
This section describes the benchmarks, architectures, scaling grid,
and training protocol in detail.

\subsection{PDE Benchmarks}
\label{sec:pdes}

We select three 1D PDEs that span a range of mathematical character:

\begin{table}[h]
\centering
\caption{PDE benchmarks. All problems map an input function
(initial condition or coefficient field) to an output function
(solution or solution at final time).}
\label{tab:tasks}
\begin{tabular}{@{}llll@{}}
\toprule
\textbf{Task} & \textbf{PDE} & \textbf{Character} &
\textbf{Solver} \\
\midrule
Burgers & $u_t + uu_x = \nu u_{xx}$ & Nonlinear, shocks &
IF-RK4 (spectral) \\
Darcy flow & $-\nabla \cdot (a \nabla u) = 1$ & Linear elliptic,
discontinuous coeff. & Spectral Galerkin \\
Diffusion & $u_t = \nu u_{xx}$ & Linear, smooth &
Exact Fourier \\
\bottomrule
\end{tabular}
\end{table}

\noindent
\textbf{Burgers equation} ($\nu = 0.01$, $T = 1.0$, periodic BC):
maps random initial conditions (superposition of Fourier modes with
$k^{-2}$ energy spectrum) to solutions at $T = 1$. The nonlinear
advection term generates shocks, requiring resolution to capture
fine-scale gradients.

\textbf{Darcy flow} (Dirichlet BC on $[0, 1]$, $f = 1$):
maps a log-normal random permeability field $a(x)$ to the pressure
solution $u(x)$. The discontinuous coefficient field makes this an
intrinsically difficult approximation problem.

\textbf{Diffusion} ($\nu = 0.01$, $T = 1.0$, periodic BC):
maps random initial conditions (same distribution as Burgers) to
solutions at $T = 1$. The heat equation smooths all fine-scale
features, so the output's spectral content is concentrated in low
Fourier modes regardless of input resolution.

These three tasks were chosen to separate three distinct resolution
regimes: resolution \emph{should} matter (Burgers, shocks),
resolution is \emph{ambiguous} (Darcy, pointwise coefficients), and
resolution \emph{should not} matter past a low threshold (diffusion,
spectral smoothing).

\subsection{Architectures}
\label{sec:architectures}

\begin{table}[h]
\centering
\caption{Architecture variants. ``Resolution dependence'' indicates
how the number of learnable parameters scales with input resolution~$R$.}
\label{tab:models}
\begin{tabular}{@{}llp{7cm}@{}}
\toprule
\textbf{Model} & \textbf{Type} & \textbf{Resolution dependence} \\
\midrule
FNO & Fourier Neural Operator & Resolution-invariant: spectral
convolutions truncate at a fixed number of modes. \\
DeepONet & Branch-Trunk network & Resolution-invariant: branch
processes $R$-dimensional input, but trunk queries at arbitrary
points. \\
MLP baseline & Feedforward MLP & $D \propto R$: input and output
layers have dimension $R$, so parameter count grows linearly with
resolution. \\
MLP controlled & Fixed-param MLP & $D \approx \mathrm{const}$:
hidden width is adjusted at each $R$ to hold total parameter count
constant (matched to $R = 64$). \\
\bottomrule
\end{tabular}
\end{table}

The MLP controlled variant is critical for disentangling the
\emph{resolution--capacity confound} in MLP-based models: when
an MLP's error decreases with resolution, is it because the model
received more informative input (a resolution effect) or because it
has more parameters (a capacity effect)?
By fixing parameter count, we isolate the resolution signal.

\subsection{Scaling Grid}
\label{sec:grid}

Each axis is varied over a range spanning approximately two orders of
magnitude:

\begin{itemize}[leftmargin=*,itemsep=2pt]
  \item \textbf{Dataset size} $N \in \{50, 100, 500, 1000, 5000\}$
    (baseline), extended to 11 values
    $N \in \{50, 100, 150, 200, 300, 500, 750, 1000, 2000, 3000, 5000\}$
    in a densified grid.
  \item \textbf{Resolution} $R \in \{32, 64, 128, 256\}$ (baseline),
    extended to 10 values
    $R \in \{16, 32, 48, 64, 96, 128, 192, 256, 384, 512\}$ for
    Burgers.
  \item \textbf{Capacity} $D$: 7 levels (tiny through xlarge) spanning
    $\sim$2K to $\sim$2.7M parameters. MLP and DeepONet share a
    hidden-width grid; FNO varies modes, width, and depth
    independently.
\end{itemize}

Every combination of $(N, R, D, \text{task}, \text{model})$ in the
grid is trained with 3 data seeds $\times$ 3 training seeds $= 9$
independent runs, for a total of \mainRuns{} main runs plus
\ablationRuns{} optimizer-ablation runs.

\begin{table}[h]
\centering
\caption{Run accounting.}
\label{tab:run-accounting}
\begin{tabular}{@{}lr@{}}
\toprule
\textbf{Category} & \textbf{Runs} \\
\midrule
Baseline grid (3 tasks $\times$ 3 models) & 15{,}120 \\
Dense-$N$ extensions & 11{,}578 \\
Wide-$R$ extensions (Burgers) & 2{,}352 \\
MLP controlled (Burgers + Diffusion) & 2{,}016 \\
Dense-$N$ MLP controlled & 304 \\
\midrule
\emph{Subtotal (main)} & \mainRuns{} \\
\midrule
Adam $2\times$ patience (Burgers) & 756 \\
SGD + cosine annealing (Burgers) & 657 \\
\midrule
\textbf{Total} & \totalRuns{} \\
\bottomrule
\end{tabular}
\end{table}

\subsection{Training Protocol}
\label{sec:training}

All models are trained with an identical protocol to ensure that
performance differences reflect architecture and scaling, not training
configuration:

\begin{itemize}[leftmargin=*,itemsep=2pt]
  \item \textbf{Optimizer:} Adam, learning rate $10^{-3}$, weight
    decay $10^{-5}$.
  \item \textbf{Convergence:} Maximum 5{,}000 epochs with early
    stopping (patience 200 epochs, monitoring validation
    relative~$L^2$).
  \item \textbf{Batch size:} 64.
  \item \textbf{Data split:} 80/10/10 train/validation/test,
    deterministic given the data seed.
  \item \textbf{Loss:} Relative $L^2$ error,
    $\relLtwo = \|\hat{v} - v\|_2 / \|v\|_2$.
  \item \textbf{Evaluation:} Test-set $\relLtwo$ at the
    best-validation checkpoint.
\end{itemize}

\noindent
Using a \emph{shared} training protocol is a deliberate choice:
it sacrifices per-architecture optimality in favor of a controlled
comparison. This design decision has a well-known tradeoff.
On one hand, it ensures that performance differences are not
confounded with hyperparameter-tuning effort. On the other hand,
if one architecture is substantially more sensitive to the optimizer,
learning-rate schedule, or architectural hyperparameters (e.g.,
DeepONet's branch/trunk width balance), a shared protocol can
turn ``controlled comparison'' into ``architecture-dependent
under-tuning.''
We partially address this concern through optimizer ablations
(\Cref{sec:ablations}) and provide an extended fairness discussion
in \Cref{sec:deeponet}.


% ══════════════════════════════════════════════════════════════════════
\section{Architecture Comparison}
\label{sec:architecture}

We begin with the most robust finding: architecture choice overwhelms
every other scaling consideration in this study.

At matched capacity (medium, $\sim$34K parameters for MLP at
$R = 128$), FNO achieves the lowest test error at \emph{every}
combination of task, dataset size, and resolution---143 out of 143
head-to-head comparisons.

\begin{table}[h]
\centering
\caption{Best achievable relative $L^2$ error per task and model
(across all $N$, $R$, and capacity levels).}
\label{tab:best-error}
\begin{tabular}{@{}lcccc@{}}
\toprule
\textbf{Task} & \textbf{FNO} & \textbf{MLP baseline} &
\textbf{DeepONet} & \textbf{MLP controlled} \\
\midrule
Burgers   & \textbf{0.027} & 0.106 & 0.115 & 0.207 \\
Darcy     & \textbf{0.471} & 0.685 & 0.743 & --- \\
Diffusion & \textbf{0.002} & 0.018 & 0.098 & 0.016 \\
\bottomrule
\end{tabular}
\end{table}

The FNO advantage is not marginal. On diffusion, FNO achieves
$\relLtwo = 0.002$ where DeepONet plateaus at 0.098---a $49\times$
gap. On Burgers, FNO outperforms the next-best model by
$4$--$6\times$ depending on the regime.

\paragraph{Interpretation.}
FNO's spectral inductive bias---learning convolutions in Fourier
space---is well matched to these periodic 1D PDEs. The strength of
the advantage suggests that for problems with smooth or spectrally
sparse solutions, the choice of inductive bias matters far more than
brute-force scaling along any axis.
We emphasize that this conclusion is specific to 1D tasks with
periodic or simple boundary conditions. Different architectures may
be competitive or superior in higher dimensions, on irregular
domains, or with more complex boundary conditions.

\paragraph{Spectral benchmark bias.}
\label{sec:spectral-bias}
The strength of FNO's dominance should be interpreted in light of an
important structural feature of our benchmark suite: two of the three
tasks---Burgers and diffusion---are periodic and solved in Fourier
space, directly matching FNO's spectral convolution kernel.
Darcy flow, while not periodic, is solved via spectral Galerkin on
a regular grid with smooth solutions, which remains amenable to
Fourier representations.
In other words, FNO's dominance is \emph{partly expected by design}:
these tasks are close to a best-case scenario for spectral methods.
For problems on irregular domains, with non-periodic boundary
conditions, or with spatially localized features (e.g., crack
propagation, fluid--structure interaction), architectures with local
inductive biases (e.g., graph neural operators, attention-based
models) may be more competitive.
The 143/143 result should therefore be read as ``FNO is dominant on
spectrally favorable 1D benchmarks,'' not as a general claim about
operator-learning architectures.


% ══════════════════════════════════════════════════════════════════════
\section{Relative Importance of Scaling Axes}
\label{sec:importance}

To provide a \emph{descriptive summary} of how much each axis
co-varies with error, we fit a log-linear regression:
\begin{equation}
\label{eq:log-linear}
\log \relLtwo = \beta_0 + \beta_N \log N + \beta_R \log R
  + \beta_D \log D + \epsilon
\end{equation}
for each task--model pair and compute the relative importance
$|\beta_i| / \sum_j |\beta_j|$.

\paragraph{Caution: associative, not causal.}
These importance percentages describe how much of the \emph{observed}
log-error variation is associated with each axis in a linear model.
They should not be interpreted as causal attributions: the axes are
varied on a designed grid (not randomized), interactions are omitted,
and the log-linear form is an approximation.
Furthermore, the $R^2$ values range from 0.35 (DeepONet on Darcy)
to 0.83 (MLP on Darcy), meaning 17--65\% of the variance is
unexplained.
The decomposition is useful as a coarse summary of which levers
matter most, but should not be treated as a precise budget
allocation tool.

\begin{table}[h]
\centering
\caption{Relative importance of each scaling axis
(percentage of total $|\beta|$), with overall $R^2$ of the
log-linear fit. $N$ dominates for all task--model pairs
except MLP controlled, where resolution dominates in the
harmful direction.}
\label{tab:importance}
\begin{tabular}{@{}llcccc@{}}
\toprule
\textbf{Task} & \textbf{Model} & $N$ & $R$ & $D$ & $R^2$ \\
\midrule
Burgers   & FNO          & \textbf{69.6\%} & 17.2\% & 13.2\% & 0.64 \\
Burgers   & DeepONet     & 43.6\% & 22.2\% & \textbf{34.2\%} & 0.49 \\
Burgers   & MLP baseline & \textbf{66.1\%} & 11.4\% & 22.5\% & 0.68 \\
Darcy     & FNO          & \textbf{85.2\%} & 8.4\%  & 6.3\%  & 0.77 \\
Darcy     & DeepONet     & \textbf{92.8\%} & 4.0\%  & 3.2\%  & 0.35 \\
Darcy     & MLP baseline & \textbf{84.9\%} & 5.1\%  & 10.0\% & 0.83 \\
Diffusion & FNO          & \textbf{82.0\%} & 0.5\%  & 17.5\% & 0.68 \\
Diffusion & MLP baseline & \textbf{82.1\%} & 13.8\% & 4.1\%  & 0.82 \\
Diffusion & MLP controlled & 12.2\% & \textbf{51.8\%} & 36.0\% & 0.46 \\
\bottomrule
\end{tabular}
\end{table}

Dataset size is associated with 65--93\% of error variance for all
standard models (in the log-linear sense described above).
Capacity is a distant second. Resolution contributes less
than 20\% in every case---and for MLP controlled on diffusion, the
51.8\% resolution importance reflects the \emph{harmful} direction
(error increases with~$R$; \Cref{sec:resolution}).

The low overall $R^2$ for DeepONet on Darcy (0.35) deserves
comment: 65\% of the error variance is unexplained by the
three scaling axes combined, reflecting the high seed-to-seed
variability on this task (\Cref{sec:darcy}).


% ══════════════════════════════════════════════════════════════════════
\section{Data Scaling: Steep but Saturating}
\label{sec:data}

\begin{figure}[t]
\centering
\includegraphics[width=\textwidth]{../figures/fig_data_scaling_per_task.pdf}
\caption{Data scaling curves (test $\relLtwo$ vs.\ $N$)
for each task and model, averaged over resolutions and capacity
levels. Error bars show $\pm 1$ standard deviation across seeds.
All models show diminishing returns past $N \approx 1000$.}
\label{fig:data-scaling}
\end{figure}

Data scaling is the most useful lever in this study. Effective
power-law exponents (slopes in the $\log N$--$\log \relLtwo$ plane,
averaged over $R$ and $D$) are:

\begin{table}[h]
\centering
\caption{Effective data-scaling exponents $\alpha_N$ (more negative
= steeper improvement with data). Averaged over all resolutions
and capacity levels.}
\label{tab:data-exponents}
\begin{tabular}{@{}lccc@{}}
\toprule
\textbf{Task} & \textbf{FNO} & \textbf{MLP baseline} &
\textbf{DeepONet} \\
\midrule
Burgers   & $-0.321$ & $-0.141$ & $-0.067$ \\
Darcy     & $-0.090$ & $-0.328$ & $-0.145$ \\
Diffusion & $-0.352$ & $-0.518$ & $-0.101$ \\
\bottomrule
\end{tabular}
\end{table}

FNO shows steep data scaling on Burgers and diffusion
($\alpha_N \approx -0.3$ to $-0.35$): a $10\times$ increase in
data roughly halves the error. DeepONet has the weakest data
scaling, with exponents near $-0.07$ to $-0.15$, indicating
early saturation of its learning capacity from additional samples.

\paragraph{Diminishing returns.}
The data--error curve bends sharply. For FNO on diffusion:
$N = 50 \to 100$ ($2\times$) reduces error by 42\%, but
$N = 1000 \to 5000$ ($5\times$) reduces error by only 0.3\%.
For FNO on Burgers: $N = 200 \to 1000$ ($5\times$) yields
$\sim$40\% improvement, but $N = 1000 \to 5000$ ($5\times$)
yields only $\sim$7\%.
Past $N \approx 1000$, additional data provides minimal benefit
for FNO across all tasks, suggesting the remaining error is
dominated by the architecture's approximation limits or the
irreducible error floor of the task at the given resolution.

\paragraph{Data efficiency.}
The practical consequence is striking: FNO on diffusion
reaches $\relLtwo < 0.02$ with just 50 samples. The spectral
structure of the heat equation perfectly matches FNO's
Fourier-space inductive bias. Burgers requires $10$--$100\times$
more data to reach the same accuracy, consistent with the
greater complexity of nonlinear shock dynamics.
Darcy does not reach $\relLtwo < 0.10$ with any
model/data/capacity combination in this study.


% ══════════════════════════════════════════════════════════════════════
\section{Resolution Is Not a Useful Scaling Axis}
\label{sec:resolution}

\begin{figure}[t]
\centering
\includegraphics[width=\textwidth]{../figures/fig_resolution_scaling_per_task.pdf}
\caption{Resolution scaling (test $\relLtwo$ vs.\ $R$) for each
task and model. For diffusion, higher resolution consistently
\emph{increases} error. For Darcy, resolution has no systematic
effect. Only for Burgers does resolution provide a benefit, and
primarily for non-spectral architectures.}
\label{fig:res-scaling}
\end{figure}

The central empirical finding of this study is that discretization
resolution is not a reliable scaling axis for neural operator
learning. The effect of resolution is task-dependent and often
harmful.

\subsection{Diffusion: Higher Resolution Increases Error}
\label{sec:res-diffusion}

On diffusion, all four model variants show monotonically increasing
error with resolution:

\begin{table}[h]
\centering
\caption{Resolution effect on diffusion (test $\relLtwo$ at matched
$N$ and capacity). Higher $R$ consistently degrades performance.}
\label{tab:res-diffusion}
\begin{tabular}{@{}lcccc@{}}
\toprule
\textbf{Model} & $R{=}32$ & $R{=}128$ & $R{=}256$ &
\textbf{Degradation} \\
\midrule
MLP baseline   & 0.082 & 0.099 & 0.098 & $+20\%$ \\
MLP controlled & 0.021 & 0.027 & 0.026 & $+24\%$ \\
DeepONet       & 0.209 & ---   & 0.243 & $+16\%$ \\
FNO            & 0.008 & ---   & 0.008 & $+12\%$ \\
\bottomrule
\end{tabular}
\end{table}

This pattern holds across \emph{every} dataset size at every
capacity level. The MLP controlled variant shows an $R$-scaling
exponent of $+0.638$ ($R^2 = 0.96$)---strong, reliable evidence
that higher resolution hurts, even when the capacity confound is
removed.

\paragraph{Mechanism.}
The heat equation smooths all fine-scale features exponentially.
At $T = 1$ with $\nu = 0.01$, Fourier modes above $k \approx 10$
are damped below machine precision. Increasing $R$ beyond $\sim$32
adds only output-irrelevant high-frequency content: the model must
learn to ignore grid points that carry no predictive information.
This wastes model capacity and, for small models, can actively
degrade predictions.

\subsection{Burgers: Resolution Helps, but with Caveats}
\label{sec:res-burgers}

For Burgers, resolution generally improves accuracy---shock structure
requires fine grids. However, the benefit is modest for FNO
(which operates in spectral space) and strongest for the MLP baseline,
where it is confounded with the capacity increase
(\Cref{sec:confound}).
The widened resolution grid (10 values, $R \in \{16, \ldots, 512\}$)
reveals that FNO resolution scaling is approximately linear in~$R$,
consistent with each additional grid point contributing one more
retained Fourier mode.

\subsection{Darcy: Resolution Is Noise}
\label{sec:res-darcy}

DeepONet $R$-exponents fluctuate between $-0.05$ and $+0.25$
depending on $N$, with per-slice $R^2 \approx 0.3$: there is no
reliable resolution signal. FNO shows $R$-exponents of $-0.03$ to
$+0.03$. Resolution changes nothing meaningful for Darcy flow.


% ══════════════════════════════════════════════════════════════════════
\section{The MLP Resolution--Capacity Confound}
\label{sec:confound}

For standard MLPs, the parameter count grows linearly with
resolution: $D = 2Rh + h^2 + 2h + R$ for hidden width $h$.
Doubling $R$ roughly doubles $D$. Any apparent ``resolution
benefit'' in MLP performance could therefore be a capacity effect.

The MLP controlled variant resolves this confound by adjusting
hidden width to hold $D$ constant across resolutions (matched to
the standard MLP at $R = 64$).

\begin{table}[h]
\centering
\caption{Diffusion: MLP controlled (fixed $D$) vs.\ MLP baseline
(growing $D$). Both show error increasing with $R$, ruling out the
capacity confound as the driver of resolution harm.}
\label{tab:mlp-controlled}
\small
\begin{tabular}{@{}lcccc@{}}
\toprule
\textbf{Capacity} & $R{=}32$ & $R{=}64$ & $R{=}128$ & $R{=}256$ \\
\midrule
\multicolumn{5}{l}{\emph{MLP controlled (fixed parameter count)}} \\
tiny      & 0.023 & 0.038 & 0.140 & 0.295 \\
small     & 0.024 & 0.023 & 0.025 & 0.111 \\
medium    & 0.021 & 0.027 & 0.027 & 0.026 \\
large     & 0.018 & 0.024 & 0.031 & 0.034 \\
xlarge    & 0.027 & 0.032 & 0.039 & 0.047 \\
\midrule
\multicolumn{5}{l}{\emph{MLP baseline (parameter count grows with $R$)}} \\
tiny      & 0.118 & 0.118 & 0.125 & 0.132 \\
small     & 0.096 & 0.101 & 0.105 & 0.103 \\
medium    & 0.082 & 0.091 & 0.099 & 0.098 \\
large     & 0.070 & 0.081 & 0.093 & 0.100 \\
xlarge    & 0.102 & 0.117 & 0.127 & 0.133 \\
\bottomrule
\end{tabular}
\end{table}

Both variants show error increasing with $R$ on diffusion.
The effect is most dramatic at small model sizes (controlled MLP
tiny: $0.023 \to 0.295$ from $R = 32$ to $R = 256$).
On Burgers, the controlled MLP shows modest resolution benefit
($\relLtwo$ drops from 0.379 to 0.277 going from $R = 32$ to
$R = 128$), confirming that resolution does carry genuine
information for shock-forming problems---but the benefit is far
smaller than the standard MLP's apparent improvement, which
conflates resolution with capacity growth.

\paragraph{Recommendation.}
When evaluating resolution effects on MLPs or other
resolution-dependent architectures, parameter count must be
controlled. Failure to do so risks attributing capacity effects
to resolution.


% ══════════════════════════════════════════════════════════════════════
\section{Capacity Scaling: Diminishing Returns}
\label{sec:capacity}

\begin{figure}[t]
\centering
\includegraphics[width=\textwidth]{../figures/fig_capacity_scaling_per_task.pdf}
\caption{Capacity scaling (test $\relLtwo$ vs.\ parameter count $D$)
for each task and model. Returns diminish sharply: 100$\times$ more
parameters often yields less than 20\% error reduction.}
\label{fig:cap-scaling}
\end{figure}

Capacity scaling exponents are uniformly weak:

\begin{table}[h]
\centering
\caption{Effective capacity-scaling exponents $\alpha_D$. All are
close to zero: massive model scaling provides minimal benefit.}
\label{tab:cap-exponents}
\begin{tabular}{@{}lccc@{}}
\toprule
\textbf{Task} & \textbf{FNO} & \textbf{MLP baseline} &
\textbf{DeepONet} \\
\midrule
Burgers   & $-0.037$ & (confounded) & $-0.066$ \\
Darcy     & $-0.004$ & $-0.063$     & $-0.003$ \\
Diffusion & $-0.023$ & $-0.014$     & $-0.107$ \\
\bottomrule
\end{tabular}
\end{table}

On Darcy, increasing FNO parameters by $388\times$ (from 7K to 2.7M)
reduces error by only 18\% ($0.573 \to 0.471$). DeepONet capacity
exponent on Darcy is $-0.003$: capacity is irrelevant.

\paragraph{Overparameterization can hurt.}
FNO on diffusion achieves its \emph{best} error at the small
capacity level (44K parameters, $\relLtwo = 0.00213$); the
xlarge model (2.7M parameters) achieves 0.00307---worse by 44\%.
When the task is simple enough to be solved with a small model,
additional capacity introduces overfitting or optimization
difficulties.

\paragraph{Contrast with language models.}
This stands in sharp contrast to the NLP scaling regime,
where capacity exponents of $-0.07$ to $-0.076$ are reported
\citep{kaplan2020scaling}---similar in magnitude but applied
across \emph{billions} of parameters, yielding large absolute
gains. In our 1D operator-learning setting, the parameter range
is much smaller ($\sim$2K--2.7M) and the tasks are comparatively
simple, so the capacity returns are exhausted early.


% ══════════════════════════════════════════════════════════════════════
\section{Regime-Dependent Scaling Forms}
\label{sec:law-form}

Do the error curves follow a universal functional form?
Following \citet{perrier2025scaling}, we fit six candidate laws
(power, exponential, logarithmic, stretched exponential, saturation,
linear) to each per-slice scaling curve and select the winner via
the corrected Akaike Information Criterion (AICc).

The principal finding of this section is \emph{methodological, not
substantive}: the identity of the AICc-selected law is fragile,
shifting with grid density, task, and architecture.
The value of the multi-law framework lies not in identifying the
``correct'' law, but in revealing when a power-law assumption is
inadequate and in diagnosing the stability of any particular fit.

\begin{figure}[t]
\centering
\includegraphics[width=\textwidth]{../figures/multilaw_heatmap.pdf}
\caption{AICc-selected law per task--model--axis combination.
No single functional form dominates; the winning law depends on the
task, architecture, and axis.}
\label{fig:heatmap}
\end{figure}

The answer is clearly \emph{no}: the winning law varies across tasks,
models, and axes. However, we stress that the inability to identify
a single winner is itself the most robust finding here.
At per-slice sample sizes of 4--7 points, even AICc
(which penalizes complexity more aggressively than AIC at small~$n$)
cannot reliably distinguish among laws whose predictions differ only
in the tails.
The regime-dependence is real, but the specific law assignments
reported below should be treated as \emph{exploratory} rather than
definitive.

\subsection{Data Axis}
On the original 5-point data grid, logarithmic scaling wins for most
task--model pairs. However, when the grid is densified to 11 points,
the winners shift systematically toward bounded-decay families
(exponential, saturation, stretched exponential)---all of which encode
a finite irreducible error floor that the logarithmic form lacks.
With only five sample sizes, the within-regime curvature is too
coarsely sampled for AICc to distinguish a log curve from a
bounded-decay curve.

\paragraph{Implication.}
Law selection on the data axis is sensitive to the number of
evaluation points. The qualitative finding---data scaling has
bounded diminishing returns---is robust, but the specific
functional form is not identifiable from typical experimental
grids.

\subsection{Resolution and Capacity Axes}
On the resolution axis (4 points), logarithmic wins for most
MLP/DeepONet slices, linear wins for FNO on Burgers. On the
widened 10-point grid, FNO clearly shifts to linear (81\% of
slices), while MLP and DeepONet retain logarithmic as a narrow
plurality. The capacity axis ($n = 7$ points) shows the most
stable law selection, likely because it has the most data points
per slice.

\subsection{Sensitivity to Grid Density: The Central Methodological Finding}
\label{sec:grid-sensitivity}
The shift in data-axis law selection from logarithmic (5 points)
to bounded-decay (11 points) is the most important methodological
finding of this section:
\emph{scaling-law identification requires dense grids, and typical
experimental designs in operator learning are too sparse to support
law-form claims}.
At $n = 4$--5 points per slice, even AICc (designed for small
samples) cannot reliably distinguish functional forms.
Within-slice bootstrap analysis (\Cref{app:bootstrap}) confirms
this fragility: only 32--64\% of data-axis slices produce a
consistent winner across bootstrap resamples, compared to 83--100\%
on the resolution axis (which shows less curvature ambiguity).
We recommend a minimum of 8--10 evaluation points per axis for any
future scaling study that aims to identify the functional form, not
merely the direction.

\paragraph{Implication for this paper's own results.}
This grid-sensitivity finding applies reflexively: our own
law-selection results on 4--5-point slices should be read as
suggestive, not definitive. The directional conclusions (data
scaling has diminishing returns, resolution can hurt) are robust
to grid density, but the specific claim that ``exponential beats
power law on the data axis'' is not.


% ══════════════════════════════════════════════════════════════════════
\section{Cross-Resolution Transfer}
\label{sec:transfer}

\begin{figure}[t]
\centering
\includegraphics[width=\textwidth]{../figures/fig5_transfer_heatmap.pdf}
\caption{Cross-resolution transfer: test $\relLtwo$ when models
trained at $\Rtrain$ are evaluated at $\Reval \neq \Rtrain$.
FNO transfers modestly upward (best at $\Reval / \Rtrain \approx
1$--$2\times$), then degrades.}
\label{fig:transfer}
\end{figure}

For resolution-invariant architectures (FNO, DeepONet), we evaluate
models trained at $\Rtrain$ on inputs at $\Reval \neq \Rtrain$.

\textbf{FNO on Burgers} performs best at $\Reval / \Rtrain \approx
1$--$2\times$, with test error of $\sim$0.25. At $16\times$ higher
evaluation resolution, error returns to 0.43---near untrained levels.
The model's learned spectral filters provide modest upward transfer
but not dramatic zero-shot super-resolution.

\textbf{FNO on Darcy} is essentially flat across all ratios
($\relLtwo \approx 0.68$--$0.70$): the task has such weak
resolution dependence that cross-resolution evaluation adds no
information.

\textbf{DeepONet on Darcy} degrades monotonically from
$\relLtwo = 1.36$ (ratio $= 0.12$) to 1.57 (ratio $= 16$),
suggesting the architecture encodes resolution-specific features
in its branch network.

An ablation comparing interpolation-based and subsampling-based
cross-resolution evaluation (\Cref{app:deeponet-ablation}) confirms
that these findings are not driven by interpolation artifacts.


% ══════════════════════════════════════════════════════════════════════
\section{Failure Modes}
\label{sec:failures}

\subsection{Darcy Flow: A Pathological Case}
\label{sec:darcy}

Darcy flow stands out as anomalous across every analysis dimension.
Because it is one of only three tasks and serves as the main
exception to several of our findings, we provide an extended
diagnosis here. Conclusions from Darcy should not be pooled
uncritically with results from Burgers and diffusion.

\subsubsection{Symptoms}

\begin{itemize}[leftmargin=*,itemsep=2pt]
  \item \textbf{High error floor.} The best achievable error is
    $\relLtwo = 0.471$ (FNO). No model reaches below 47\% relative
    error---compared to 2.7\% (Burgers) and 0.2\% (diffusion).

  \item \textbf{Extreme variance.} DeepONet on Darcy has a
    coefficient of variation (CV) of up to 1.54 across seeds. Mean
    error values of $\relLtwo > 1$ (worse than predicting the mean)
    appear frequently.

  \item \textbf{Pervasive non-monotonicity.} Dozens of configurations
    show \emph{increasing} error when data, resolution, or capacity
    is added. DeepONet at multiple capacity levels shows
    $\relLtwo$ jumping from 0.80 to 1.88 when going from $N = 1000$
    to $N = 2000$.

  \item \textbf{Low explanatory power.} The log-linear regression
    $R^2$ for DeepONet on Darcy is only 0.35: none of the three
    scaling axes adequately explains the variance.

  \item \textbf{Trivial-baseline failures.} Multiple DeepONet
    configurations produce $\relLtwo > 1.0$, meaning they perform
    worse than a trivial predictor that outputs the training-set
    mean. This occurs at 12\% of DeepONet/Darcy configurations
    (large and xlarge capacity, various $N$ and $R$).
\end{itemize}

\subsubsection{Diagnosis}

We identify three likely contributing factors:

\paragraph{1. Discontinuous coefficients break spectral and
  pointwise approximation.}
The Darcy permeability field $a(x)$ is drawn from a log-normal
random field and can exhibit near-discontinuities. Unlike Burgers
and diffusion, where both input and output functions are smooth
(or at least piecewise smooth with known structure), the Darcy
input--output map requires learning a nonlinear relationship between
a rough coefficient field and a smooth solution.
This is a fundamentally harder approximation class: the
Kolmogorov $n$-width of the Darcy solution manifold under
coefficient perturbations decays more slowly than for Burgers
or diffusion, implying that all linear methods (including FNO's
spectral truncation) face a harder problem.

\paragraph{2. The task is not periodic.}
Darcy uses Dirichlet boundary conditions on $[0, 1]$. FNO's
periodic Fourier convolutions introduce boundary artifacts
(Gibbs-like oscillations) when applied to non-periodic data,
which may partly explain why FNO's Darcy error floor (0.471) is
$170\times$ higher than its diffusion error floor (0.002).
This is consistent with the spectral-bias interpretation from
\Cref{sec:spectral-bias}: when the problem structure does not
align with the spectral inductive bias, the advantage shrinks
dramatically.

\paragraph{3. The optimization landscape is rugged for DeepONet.}
The extreme seed-to-seed variance (CV up to 1.54) and
non-monotonic scaling curves suggest that DeepONet's
branch--trunk decomposition encounters optimization pathologies
on this task.
The branch network must encode the relationship between a
rough ($a(x)$) input and the trunk's smooth basis, creating a
loss landscape with many local minima.
This interpretation is supported by the observation that
non-monotonicity concentrates at larger capacity levels
(large/xlarge), where the optimization problem is harder.

\subsubsection{Impact on paper conclusions}

Because Darcy is pathological, several aggregate statistics reported
elsewhere in the paper are influenced by its inclusion:
\begin{itemize}[leftmargin=*,itemsep=2pt]
  \item The log-linear importance percentages for DeepONet
    (\Cref{tab:importance}) have low reliability on Darcy ($R^2 = 0.35$).
  \item Capacity-scaling exponents for Darcy ($\alpha_D \approx 0$)
    may reflect optimization failure rather than genuine capacity
    saturation.
  \item Resolution-axis conclusions (``resolution is noise for Darcy'')
    are confounded with the high variance: the signal may exist but
    be undetectable.
\end{itemize}
We therefore recommend that readers treat Darcy results as a
separate case study rather than pooling them with Burgers and
diffusion when drawing scaling-law conclusions.

\subsection{DeepONet Saturation and the Shared-Protocol Fairness Question}
\label{sec:deeponet}

DeepONet exhibits two distinct failure modes, but interpreting them
requires confronting a methodological concern: the shared training
protocol may systematically disadvantage DeepONet relative to FNO
and MLP.

\subsubsection{The fairness concern}

Our protocol fixes Adam with learning rate $10^{-3}$, weight decay
$10^{-5}$, and patience 200 for all architectures. This ensures a
controlled comparison, but it creates a potential confound:
if DeepONet is more sensitive than FNO or MLP to learning-rate
schedule, optimizer choice, or branch/trunk width balance, then
the shared protocol effectively \emph{under-tunes} DeepONet while
appropriately tuning FNO.

Several observations suggest this concern is real but does
not fully explain the performance gap:

\begin{enumerate}[leftmargin=*,itemsep=2pt]
  \item \textbf{Optimizer ablation evidence.}
    On Burgers, switching from Adam to SGD+cosine for DeepONet
    changes the data-scaling exponent (from power to exponential
    winner family) but does not close the gap to FNO: DeepONet's
    best error under SGD+cosine remains $3$--$5\times$ worse than
    FNO under the baseline Adam protocol.
    Similarly, doubling early-stopping patience from 200 to 400
    epochs does not improve DeepONet's ranking.

  \item \textbf{DeepONet-specific tuning is known to help.}
    The literature reports that DeepONet benefits from
    branch/trunk learning-rate balancing~\citep{wang2021learning},
    which our protocol does not include.
    We cannot rule out that architecture-specific tuning would
    narrow the gap, and this is an important limitation.

  \item \textbf{The saturation pattern is architecture-specific.}
    DeepONet's early saturation (data exponent $\alpha_N \approx
    -0.04$ to $-0.07$ on Burgers at $R \geq 32$) is an order of
    magnitude weaker than FNO's ($\alpha_N \approx -0.32$).
    If this were purely an optimizer issue, we would expect the
    gap to be approximately constant across $N$; instead, it
    \emph{widens} with increasing data, consistent with an
    expressiveness bottleneck in the branch--trunk decomposition
    rather than an optimization failure.

  \item \textbf{MLP under the same protocol performs better.}
    On Burgers and diffusion, the MLP baseline outperforms DeepONet
    under the identical protocol. Since MLP has no special tuning
    advantage, DeepONet's relative weakness is not solely
    attributable to under-tuning.
\end{enumerate}

\paragraph{Assessment.}
The shared protocol likely under-optimizes DeepONet to some degree.
However, the FNO--DeepONet gap ($4$--$49\times$ depending on task) is
too large to be plausibly explained by optimizer choice alone given that
the ablation already tests two alternative optimizers with minimal
change in the ranking.
We conclude that the qualitative ranking (FNO $\gg$ MLP $\gtrsim$
DeepONet on these tasks) is robust, but the \emph{magnitude} of
DeepONet's disadvantage is likely overstated by the shared protocol.
A full architecture-specific hyperparameter search would strengthen
this comparison, and we flag this as the most important follow-up
experiment.

\subsubsection{Failure modes}

\paragraph{Early saturation.}
On Burgers, DeepONet's data-scaling exponent ($\alpha_N \approx
-0.04$ to $-0.07$ at $R \geq 32$) is an order of magnitude
weaker than FNO's. Even with a $100\times$ increase in data
($N = 50 \to 5000$), error drops only $\sim$25\%, compared to FNO's
$\sim$85\% reduction. On diffusion, DeepONet plateaus at
$\relLtwo \sim 0.10$--$0.22$ regardless of data, while FNO
reaches 0.002.

\paragraph{More data can hurt.}
On Burgers, DeepONet with large/xlarge capacity consistently
shows 5--15\% error \emph{increases} when going from $N = 2000$
to $N = 5000$, across multiple resolutions. This may reflect
the branch-trunk decomposition reaching expressiveness saturation
before the training set is fully utilized, leading to degraded
optimization at larger~$N$.


% ══════════════════════════════════════════════════════════════════════
\section{Reproducibility and Noise}
\label{sec:noise}

All configurations are replicated across 9 seeds
(3~data $\times$ 3~training). \Cref{tab:noise} reports the
coefficient of variation (CV $= \sigma / \mu$) across seeds:

\begin{table}[h]
\centering
\caption{Mean coefficient of variation (CV) across seeds per
task--model pair. Darcy is $3$--$8\times$ noisier than other tasks.}
\label{tab:noise}
\begin{tabular}{@{}lccc@{}}
\toprule
\textbf{Task} & \textbf{FNO} & \textbf{MLP baseline} &
\textbf{DeepONet} \\
\midrule
Burgers   & 0.08 & 0.03 & 0.05 \\
Darcy     & 0.07 & 0.25 & 0.25 \\
Diffusion & 0.17 & 0.07 & 0.11 \\
\bottomrule
\end{tabular}
\end{table}

Darcy is extremely noisy for non-FNO architectures (CV $= 0.25$),
meaning a typical run's error varies by $\pm 25\%$ of the mean
from seed changes alone. The noisiest individual configuration
(DeepONet/med-large/$N = 2000$/$R = 256$ on Darcy, CV $= 1.54$)
produces runs ranging from near-zero to several times the mean
error. This noise level severely limits the reliability of
single-seed evaluations and underscores the importance of the
multi-seed design.


% ══════════════════════════════════════════════════════════════════════
\section{Robustness and Ablations}
\label{sec:ablations}

\subsection{Optimizer Sensitivity and Protocol Fairness}
\label{sec:optimizer}

All main results use Adam with learning rate $10^{-3}$ and patience
200. To test whether architectural rankings and directional findings
are robust to training configuration, we re-run Burgers (all three
models, all capacities) under two alternative configurations:
\begin{itemize}[leftmargin=*,itemsep=2pt]
  \item \textbf{Variant A:} Adam with $2\times$ patience (400 epochs)
  \item \textbf{Variant B:} SGD with momentum 0.9 + cosine annealing
\end{itemize}
for a total of 1{,}413 ablation runs.

\textbf{Architectural ranking:} FNO retains the best performance
under all three training configurations. The FNO--DeepONet gap
narrows slightly under SGD+cosine but remains $>3\times$ at all
capacity levels. This provides partial evidence that the ranking
is not solely a protocol artifact, though we acknowledge the
ablation is limited to Burgers and to two alternative configurations
(not a full hyperparameter search).

\textbf{Resolution axis:} Law selection is fully stable.
Logarithmic wins in 88\% of baseline slices, 90\% for Adam
$2\times$, and 100\% for SGD cosine. This is the strongest
invariance result in the study.

\textbf{Data axis:} The specific law identity is more sensitive
(DeepONet shifts from power to exponential; MLP shifts from
saturation to power), but all winners belong to bounded-decay
families. The qualitative conclusion---data scaling is bounded,
not unbounded---is optimizer-invariant.

\textbf{Conclusion:} Directional findings (which axis helps, which
hurts) are robust to the training configuration. Architectural
rankings are stable in ordering but the magnitude of gaps may be
protocol-dependent, especially for DeepONet (\Cref{sec:deeponet}).
The specific functional form on the data axis is training-sensitive,
consistent with the grid-density sensitivity discussed in
\Cref{sec:grid-sensitivity}.

\subsection{Densified Data Grid}
\label{sec:dense-n}

Extending the data axis from 5 to 11 values
($N \in \{50, \ldots, 5000\}$ with 6 intermediate points) produces
two systematic changes: (1) logarithmic is displaced as the dominant
winner in favor of bounded-decay families (exponential, saturation);
(2) AICc-BIC concordance improves from $\sim$10\% to $>$50\%.
Both are consistent with the denser grid resolving the curvature
near the error floor that was ambiguous at 5 points.

\subsection{Widened Resolution Grid}
\label{sec:wide-r}

Extending the Burgers resolution grid from 4 to 10 values
($R \in \{16, \ldots, 512\}$) reveals that FNO resolution scaling
is clearly linear (81\% of slices), while MLP and DeepONet retain
logarithmic as a narrow plurality with significant linear
competition. The directional finding---resolution helps for
Burgers---is stable regardless of grid density.


% ══════════════════════════════════════════════════════════════════════
\section{Discussion}
\label{sec:discussion}

\subsection{Summary of Evidence}

\Cref{tab:evidence} classifies findings by evidence strength.
We distinguish between findings that are robust within the 1D setting
and those that might generalize, and explicitly flag the
descriptive nature of regression-based results.

\begin{table}[h]
\centering
\caption{Summary of findings and evidence strength. All findings
are restricted to the 1D PDE setting studied here. ``Descriptive''
indicates results based on associative models (log-linear
regression) rather than controlled experiments.}
\label{tab:evidence}
\small
\begin{tabular}{@{}p{6.5cm}cp{5cm}@{}}
\toprule
\textbf{Finding} & \textbf{Strength} &
\textbf{Practical implication} \\
\midrule
FNO dominates all comparisons (143/143) on spectrally favorable 1D benchmarks &
Very strong &
Use FNO for periodic 1D operator learning \\
Dataset size co-varies most with error (65--93\%; descriptive) &
Strong &
Prioritize data collection over model tuning \\
Resolution hurts on smooth problems &
Strong &
Use the lowest $R$ that captures the physics \\
Capacity has weak returns ($\alpha_D \approx 0$) &
Strong &
A modest model suffices for 1D tasks \\
Scaling-law identification is unstable under sparse grids &
Strong &
Use $\geq$8--10 points per axis; do not assume power law \\
DeepONet saturates early (partly protocol-dependent) &
Moderate--Strong &
Consider alternatives; tune before concluding \\
Darcy is not well-served by any model &
Strong &
Requires architectural innovation \\
Cross-resolution transfer is modest &
Moderate &
Do not rely on low-$R$ training for high-$R$
deployment \\
\bottomrule
\end{tabular}
\end{table}

\subsection{Benchmark Selection Bias}
\label{sec:benchmark-bias}

A central limitation of this study is that the benchmark suite is
structurally biased toward spectral methods.
Two of three tasks (Burgers, diffusion) are periodic, solved in
Fourier space, and produce outputs whose spectral content either
decays rapidly (diffusion) or is concentrated in a known frequency
band (Burgers at moderate viscosity).
Darcy, while not periodic, uses a regular grid and produces smooth
solutions, which are still amenable to Fourier representation.

This means that our strongest finding---FNO's 143/143 dominance---is
partly a statement about benchmark--architecture alignment rather
than a statement about operator learning in general.
To contextualize the strength of this claim, consider that:
\begin{itemize}[leftmargin=*,itemsep=2pt]
  \item FNO's advantage on Darcy (the least spectrally favorable task)
    is $0.471$ vs.\ $0.685$ ($1.5\times$), far smaller than on
    diffusion ($0.002$ vs.\ $0.098$, $49\times$).
  \item The Darcy error floor of $0.471$ for FNO suggests that even
    the best-performing architecture has fundamental limitations on
    non-periodic tasks.
  \item No graph neural operators, attention-based operators, or
    other architectures with local inductive biases are included.
\end{itemize}

A stronger version of this study would include 2D/3D tasks,
non-periodic domains, multi-physics couplings, and architectures like
GNOT~\citep{hao2023gnot} or graph neural operators. Until such
evidence is available, FNO's dominance should be understood as
conditional on spectrally favorable problem structure.

\subsection{Resolution as Physics vs.\ Resolution as Dimensionality}
\label{sec:res-discussion}

The resolution axis conflates two distinct effects that are often
treated as one:

\begin{enumerate}[leftmargin=*,itemsep=2pt]
  \item \textbf{Resolution as physics information.}
    Higher resolution captures finer spatial features (shocks,
    discontinuities). For Burgers, where the solution develops
    sharp gradients, additional grid points carry genuinely new
    information.

  \item \textbf{Resolution as input dimensionality.}
    Higher resolution increases the size of the input/output vectors.
    For diffusion, where the solution is spectrally sparse, additional
    grid points beyond $R \approx 32$ carry no new information but
    increase the dimensionality of the optimization problem.
\end{enumerate}

Our results show that the second effect can dominate: when the
additional grid points carry noise or irrelevant high-frequency
content, models (especially small models) are better off at lower
resolution.
This has direct implications for benchmark design: reporting
operator-learning accuracy only at high resolution conflates learning
ability with the model's tolerance for input dimensionality.

\subsection{Implications for Practice}

For practitioners working with 1D PDEs similar to our benchmarks
(periodic, spectrally structured, on regular grids):
\begin{enumerate}[leftmargin=*,itemsep=2pt]
  \item \textbf{Architecture first.} On spectrally favorable problems,
    the choice of inductive bias (spectral vs.\ pointwise vs.\ MLP)
    has a larger effect than any amount of data, resolution, or
    capacity scaling. On non-periodic or spatially heterogeneous
    problems, this advantage may be smaller (\Cref{sec:benchmark-bias}).
  \item \textbf{Data second.} Collecting $\sim$500--1000 high-quality
    samples provides most of the achievable accuracy. Returns diminish
    sharply beyond this range.
  \item \textbf{Resolution carefully.} Match resolution to the problem's
    spectral content, not to the solver's capability. Using the
    highest available resolution can actively hurt.
  \item \textbf{Capacity modestly.} A medium-sized model ($\sim$30--50K
    parameters) captures most of the available accuracy for these 1D
    tasks; scaling to millions of parameters provides negligible
    returns.
  \item \textbf{Tune per architecture.} Our shared-protocol results
    likely understate the performance of architectures (like DeepONet)
    that benefit from specialized tuning. Before discarding an
    architecture based on benchmark comparisons, verify that the
    training protocol is not unfairly limiting it.
\end{enumerate}

\subsection{Comparison with NLP Scaling Laws}
\label{sec:nlp-comparison}

The NLP scaling regime~\citep{kaplan2020scaling,hoffmann2022training}
features three properties that our 1D PDE setting does not share:
(1)~power-law relationships spanning many orders of magnitude in
model size; (2)~weak diminishing returns across the entire tested
range; (3)~the same functional form (power law) governing all axes.
In our setting, (1)~the parameter range is modest
($\sim$2K--2.7M), (2)~diminishing returns set in by $N \approx 1000$
and $D \approx$ small, and (3)~the functional form is
regime-dependent (and fragile under bootstrap resampling,
\Cref{app:bootstrap}).
This comparison should not be over-interpreted: the NLP scaling regime
involves billions of parameters, trillions of tokens, and tasks of
far greater complexity. Our 1D PDE setting may simply be too small
and too simple for universal scaling behavior to emerge. Whether
scaling operator learning to higher dimensions, larger models, and
more complex PDEs would recover NLP-like power laws is an open
question that this study cannot answer.

\subsection{Limitations}
\label{sec:limitations}

\begin{enumerate}[leftmargin=*,itemsep=2pt]
  \item \textbf{1D only, spectrally biased benchmarks.} All tasks are
    1D PDEs, and two of three are periodic and spectrally structured.
    Resolution may be far more important in 2D/3D, where
    discretization error compounds with dimensionality.
    FNO's dominance is partly expected given this benchmark structure
    (\Cref{sec:benchmark-bias}).
    We make no claims about higher-dimensional, non-periodic, or
    irregular-domain settings.

  \item \textbf{Fixed training protocol may under-tune DeepONet.}
    A shared optimizer and schedule ensures controlled comparison but
    may systematically disadvantage architectures that are more
    sensitive to training configuration.
    The optimizer ablation partially addresses this but tests only
    two alternatives on one task---not a full hyperparameter search.
    The qualitative ranking (FNO $\gg$ others) is likely robust;
    the quantitative gap magnitude is less certain
    (\Cref{sec:deeponet}).

  \item \textbf{Darcy is pathological and destabilizes aggregate
    statistics.}
    One of three tasks produces error floors above 47\%,
    extreme seed variance, and non-monotonic scaling curves.
    Aggregate conclusions that include Darcy should be treated with
    caution (\Cref{sec:darcy}).

  \item \textbf{Log-linear importance decomposition is descriptive.}
    The relative-importance percentages (\Cref{tab:importance}) are
    associative summaries from a linear model with modest $R^2$
    values, not causal attributions (\Cref{sec:importance}).

  \item \textbf{Scaling-law forms are exploratory at sparse grids.}
    At 4--5 points per slice, law identification is fragile.
    The directional conclusions (which axis helps/hurts) are robust;
    the specific functional forms are not (\Cref{sec:grid-sensitivity}).

  \item \textbf{Limited resolution range for some tasks.}
    Darcy uses only 4~resolution values in the baseline grid.
    The wide-$R$ extension is complete only for Burgers.

  \item \textbf{Relative $L^2$ only.}
    Different error norms (pointwise max, Sobolev, spectral)
    could change accuracy rankings, especially for Burgers
    (localized shocks).

  \item \textbf{Canonical PDEs only.}
    Burgers, Darcy, and diffusion are standard benchmarks, not
    representative of the full complexity of scientific applications.
    Scaling behavior on turbulence, multi-physics, or real-world
    engineering problems may differ substantially.

  \item \textbf{No architecture-diversity beyond FNO/DeepONet/MLP.}
    Graph neural operators, transformer-based operators (GNOT),
    and physics-informed variants are not tested. Including these
    would strengthen the architecture-comparison claims.
\end{enumerate}


% ══════════════════════════════════════════════════════════════════════
\section{Conclusion}
\label{sec:conclusion}

We have presented a controlled empirical study of scaling laws for
neural operator learning across three canonical 1D PDE tasks,
three architectures, and two orders of magnitude along dataset size,
resolution, and model capacity, totaling \totalRuns{} training runs.

The study yields three central messages, all specific to the 1D
setting and benchmark suite studied here.
First, for 1D problems with periodic or simple boundary conditions,
architecture choice (specifically, FNO's spectral inductive bias)
dominates all other scaling considerations by a wide margin---though
this dominance is partly explained by the spectral structure of our
benchmark tasks, and should not be extrapolated to non-periodic or
irregular-domain settings without further evidence
(\Cref{sec:benchmark-bias}).
Second, dataset size is the only reliable scaling axis
(in the descriptive, associative sense of the log-linear analysis):
resolution is task-dependent (helpful for shocks, harmful for smooth
solutions, irrelevant for elliptic problems) and capacity returns
are negligible past modest model sizes.
Third, scaling-law identification is the most fragile aspect of our
analysis: the AICc-selected functional form shifts with grid density,
and no single law is robust across tasks and architectures.
The strongest methodological conclusion is that law identification
at typical experimental grid densities (4--5 points) is unreliable,
and future studies should use at least 8--10 evaluation points per
axis.

We also highlight two important caveats. Darcy flow behaves
pathologically across all metrics and should be treated as a separate
case study rather than pooled with Burgers and diffusion
(\Cref{sec:darcy}).
The shared training protocol may understate DeepONet's capabilities,
and the magnitude (though likely not the direction) of the
FNO--DeepONet gap should be interpreted with this in mind
(\Cref{sec:deeponet}).

These results suggest that the era of ``scaling everything'' in
operator learning may not mirror NLP's experience, at least for
1D benchmarks.
Practitioners should prioritize choosing the right
architecture for the problem's mathematical structure and
collecting sufficient training data, over scaling model size or
discretization resolution.
The extent to which these conclusions transfer to higher-dimensional,
multi-physics, or industrial-scale problems---and to architectures
with non-spectral inductive biases---remains the most important
open question.


% ══════════════════════════════════════════════════════════════════════
\begin{thebibliography}{99}

\bibitem[Chen and Chen(1995)]{chen1995universal}
T.~Chen and H.~Chen.
\newblock Universal approximation to nonlinear operators by neural networks
  with arbitrary activation functions and its application to dynamical systems.
\newblock \emph{IEEE Trans. on Neural Networks}, 6(4):911--917, 1995.

\bibitem[Hao et~al.(2023)]{hao2023gnot}
Z.~Hao, Z.~Wang, H.~Su, C.~Ying, Y.~Dong, S.~Liu, Z.~Cheng, J.~Song, and
  J.~Zhu.
\newblock {GNOT}: A general neural operator transformer for operator learning.
\newblock In \emph{ICML}, 2023.

\bibitem[Hestness et~al.(2017)]{hestness2017deep}
J.~Hestness, S.~Narang, N.~Ardalani, G.~Diamos, H.~Jun, H.~Kianinejad,
  M.~M.~A. Patwary, Y.~Yang, and Y.~Zhou.
\newblock Deep learning scaling is predictable, empirically.
\newblock \emph{arXiv:1712.00409}, 2017.

\bibitem[Hoffmann et~al.(2022)]{hoffmann2022training}
J.~Hoffmann, S.~Borgeaud, A.~Mensch, E.~Buchatskaya, T.~Cai, E.~Rutherford,
  D.~de~Las~Casas, L.~A. Hendricks, J.~Welbl, A.~Clark, T.~Hennigan,
  E.~Noland, K.~Millican, G.~van~den~Driessche, B.~Damoc, A.~Guy,
  S.~Osindero, K.~Simonyan, E.~Elsen, J.~W. Rae, O.~Vinyals, and
  L.~Sifre.
\newblock Training compute-optimal large language models.
\newblock In \emph{NeurIPS}, 2022.

\bibitem[Kaplan et~al.(2020)]{kaplan2020scaling}
J.~Kaplan, S.~McCandlish, T.~Henighan, T.~P. Brown, B.~Chess, R.~Child,
  S.~Gray, A.~Radford, J.~Wu, and D.~Amodei.
\newblock Scaling laws for neural language models.
\newblock \emph{arXiv:2001.08361}, 2020.

\bibitem[Li et~al.(2021)]{li2021fourier}
Z.~Li, N.~Kovachki, K.~Azizzadenesheli, B.~Liu, K.~Bhatt, A.~Stuart, and
  A.~Anandkumar.
\newblock Fourier neural operator for parametric partial differential equations.
\newblock In \emph{ICLR}, 2021.

\bibitem[Lu et~al.(2021)]{lu2021learning}
L.~Lu, P.~Jin, G.~Pang, Z.~Zhang, and G.~E. Karniadkis.
\newblock Learning nonlinear operators via {DeepONet} based on the universal
  approximation theorem of operators.
\newblock \emph{Nature Machine Intelligence}, 3(3):218--229, 2021.

\bibitem[Perrier(2025)]{perrier2025scaling}
L.~Perrier.
\newblock Scaling-law diagnostics for physics-informed machine learning.
\newblock Working paper, 2025.

\bibitem[Wang et~al.(2021)]{wang2021learning}
S.~Wang, H.~Wang, and P.~Perdikaris.
\newblock Learning the solution operator of parametric partial differential
  equations with physics-informed {DeepONets}.
\newblock \emph{Science Advances}, 7(40), 2021.

\bibitem[Wen et~al.(2022)]{wen2022ufno}
G.~Wen, Z.~Li, K.~Azizzadenesheli, A.~Anandkumar, and S.~M. Benson.
\newblock {U-FNO}---an enhanced {Fourier} neural operator-based deep-learning
  model for multiphase flow.
\newblock \emph{Advances in Water Resources}, 163:104180, 2022.

\end{thebibliography}


% ══════════════════════════════════════════════════════════════════════
\appendix

\section{Capacity Grid}
\label{app:capacity}

\begin{table}[h]
\centering
\caption{MLP and DeepONet capacity grid. FNO uses a separate grid
over modes, width, and number of layers.}
\label{tab:capacity-grid}
\begin{tabular}{@{}lrrr@{}}
\toprule
\textbf{Level} & \textbf{Layers} & \textbf{Width} &
$\sim D$ \textbf{(MLP, $R{=}128$)} \\
\midrule
tiny      & 2 & 32  & $\sim$5K \\
small     & 2 & 64  & $\sim$13K \\
small-med & 2 & 96  & $\sim$22K \\
medium    & 2 & 128 & $\sim$34K \\
med-large & 2 & 192 & $\sim$62K \\
large     & 2 & 256 & $\sim$82K \\
xlarge    & 3 & 256 & $\sim$148K \\
\bottomrule
\end{tabular}
\end{table}


\section{PDE Solver Details}
\label{app:solvers}

\paragraph{Burgers equation.}
$u_t + u u_x = \nu u_{xx}$ on $[0, 2\pi)$ with periodic boundary
conditions, $\nu = 0.01$, $T = 1.0$. Solved with integrating-factor
RK4 in Fourier space with $\frac{2}{3}$-rule dealiasing and
$\Delta t = 10^{-3}$. Random initial conditions are superpositions of
Fourier modes with random phases and amplitudes drawn from a $k^{-2}$
energy spectrum.

\paragraph{Darcy flow.}
$-\nabla \cdot (a(x) \nabla u(x)) = f(x)$ on $[0, 1]$ with Dirichlet
conditions, $f(x) = 1$. The coefficient field $a(x)$ is drawn from
a log-normal random field. Solved via spectral Galerkin method.

\paragraph{Diffusion equation.}
$u_t = \nu u_{xx}$ on $[0, 2\pi)$ with periodic boundary conditions,
$\nu = 0.01$, $T = 1.0$. Exact spectral solution:
$\hat{u}_k(t) = \hat{u}_k(0) e^{-\nu k^2 t}$.
Same initial-condition distribution as Burgers.


\section{Multi-Law Fitting Details}
\label{app:fitting}

Six candidate laws are fit via \texttt{scipy.optimize.curve\_fit}
(Levenberg--Marquardt) with maximum 10{,}000 iterations:

\begin{table}[h]
\centering
\caption{Candidate functional forms. $X$ is the scaling variable
($N$, $R$, or $D$); $\Einf$ denotes the irreducible error floor.}
\label{tab:candidates}
\begin{tabular}{@{}llc@{}}
\toprule
\textbf{Family} & \textbf{Form} & \textbf{Parameters} \\
\midrule
Power       & $\Einf + a X^{-\alpha}$   & 3 \\
Exponential & $\Einf + a e^{-\alpha X}$ & 3 \\
Logarithmic & $a - b \ln X$             & 2 \\
Stretched exp. & $\Einf + a e^{-(\alpha X)^\beta}$ & 4 \\
Saturation  & $\Einf + a / (X + b)$     & 3 \\
Linear      & $a + b X$                 & 2 \\
\bottomrule
\end{tabular}
\end{table}

\noindent
AICc is computed as
$\mathrm{AICc} = n \ln(\mathrm{RSS}/n) + 2k + 2k(k+1)/(n-k-1)$,
where $n$ is the number of points per slice and $k$ is the number of
parameters. When $n \leq k + 1$, AICc is set to $+\infty$.


\section{DeepONet Cross-Resolution Ablation}
\label{app:deeponet-ablation}

We compared interpolation-based and subsampling-based cross-resolution
evaluation across 288 comparisons (Darcy and Diffusion, 3 capacities
$\times$ 4 training resolutions $\times$ 4--5 evaluation resolutions
$\times$ 4 dataset sizes).

When $\Reval > \Rtrain$ (upsampling): interpolation and subsampling
perform comparably (mean $\relLtwo$: 0.523 vs.\ 0.520), with
subsampling winning 68\% of 180 comparisons by a slim margin.
When $\Reval < \Rtrain$ (downsampling): interpolation is overwhelmingly
better (98\% of 108 cases; mean $\relLtwo$: 0.531 vs.\ 1.068).

The standard interpolation strategy is robust: cross-resolution
transfer conclusions are not driven by interpolation artifacts.


\section{BIC Concordance}
\label{app:bic}

Because AICc and BIC penalties diverge at small sample sizes
(at $n = 5$, BIC's penalty $k \ln n \approx 1.6k$ is \emph{weaker}
than AICc's corrected penalty), low concordance at small $n$ reflects
the well-understood divergence between criteria, not evidence against
the AICc-selected law.

\begin{table}[h]
\centering
\caption{AICc--BIC concordance rate per axis.}
\label{tab:bic}
\begin{tabular}{@{}llccc@{}}
\toprule
\textbf{Task} & \textbf{Model} & \textbf{Data} & \textbf{Resolution} &
\textbf{Capacity} \\
\midrule
Burgers   & MLP      & 4\%  & 3\%   & 90\% \\
Burgers   & DeepONet & 21\% & 3\%   & 85\% \\
Burgers   & FNO      & 11\% & 9\%   & 70\% \\
Darcy     & MLP      & 4\%  & 43\%  & 80\% \\
Darcy     & DeepONet & 11\% & 69\%  & 95\% \\
Darcy     & FNO      & 14\% & 34\%  & 65\% \\
Diffusion & MLP      & 4\%  & 97\%  & 45\% \\
Diffusion & DeepONet & 11\% & 100\% & 80\% \\
Diffusion & FNO      & 21\% & 74\%  & 90\% \\
\bottomrule
\end{tabular}
\end{table}


\section{Within-Slice Bootstrap Stability}
\label{app:bootstrap}

To quantify the fragility of law selection, we perform a within-slice
bootstrap analysis: for each scaling slice (fixed task, model,
capacity/resolution, varying one axis), we resample the seed-averaged
data points with replacement 1{,}000 times and re-run AICc model
selection on each resample.
The ``stability'' of a slice is the fraction of bootstrap resamples
that select the same winner as the original data.

\begin{table}[h]
\centering
\caption{Within-slice bootstrap stability of AICc law selection.
Values report the fraction of slices where the original winner is
recovered in $>$50\% of bootstrap resamples. The data axis is
substantially less stable than the resolution or capacity axes.}
\label{tab:bootstrap}
\begin{tabular}{@{}llccc@{}}
\toprule
\textbf{Task} & \textbf{Model} & \textbf{Data} &
\textbf{Resolution} & \textbf{Capacity} \\
\midrule
Burgers   & MLP      & 32\%  & 97\%  & 86\% \\
Burgers   & DeepONet & 61\%  & 100\% & 71\% \\
Burgers   & FNO      & 64\%  & 83\%  & 79\% \\
Darcy     & MLP      & 43\%  & 89\%  & 80\% \\
Darcy     & DeepONet & 36\%  & 94\%  & 65\% \\
Darcy     & FNO      & 57\%  & 91\%  & 75\% \\
Diffusion & MLP      & 39\%  & 100\% & 70\% \\
Diffusion & DeepONet & 46\%  & 100\% & 83\% \\
Diffusion & FNO      & 54\%  & 97\%  & 88\% \\
\bottomrule
\end{tabular}
\end{table}

\noindent
The data axis shows stability rates of 32--64\%, meaning that for
roughly half of all data-scaling slices, changing the random seed
would change the AICc-selected law.
In contrast, resolution-axis stability is 83--100\% and capacity-axis
stability is 65--88\%.
This quantitatively confirms the methodological warning in
\Cref{sec:grid-sensitivity}: data-axis law selection at typical
sample sizes is too fragile for definitive claims about functional
form.


\end{document}
